\documentclass[10pt,a4paper]{amsart}
\usepackage[margin=19mm]{geometry}
\usepackage{amsmath,amssymb,mathtools}
\usepackage[scr=rsfs,cal=euler]{mathalfa}
\usepackage{xfrac}
\usepackage[unicode,hidelinks,bookmarks=false]{hyperref}
\newcommand{\ie}{i.e.\ }
\newcommand{\cf}{cf.\ }
\newcommand{\resp}{respectively\ }
\newcommand{\Subsection}{\subsection}
\providecommand{\nobreakdash}{\nobreak\hbox{-}}
\setlength{\emergencystretch}{2em}
\multlinegap0pt
\usepackage{bm}
\usepackage{bbm}
\usepackage{dsfont}
\usepackage{xspace}
\usepackage{cancel}
\usepackage{mathtools}
\usepackage{amsthm}
\makeatletter
\@ifundefined{theorem}{\newtheorem{theorem}{Theorem}}{}
\@ifundefined{lemma}{\newtheorem{lemma}[theorem]{Lemma}}{}
\makeatother
\newcommand{\Ad}{\operatorname{Ad}}
\newcommand{\F}{\mathbb{F}}
\newcommand{\id}{\operatorname{id}}
\title[Module-Valued 2-Local Derivations on Reductive Lie Algebras]{Module-Valued 2-Local Derivations on Reductive Lie Algebras}
\author{Yang Chen, Yongqi Luo, Junzi Xu}
\date{Source: arXiv:2608.20681v1 (CC BY 4.0)}
\begin{document}
\maketitle

\noindent\emph{Source and licence.} Module-Valued 2-Local Derivations on Reductive Lie Algebras. Version arXiv:2608.20681v1 (CC BY 4.0), distributed under the Creative Commons Attribution 4.0 International licence, as recorded in the arXiv licence metadata for this exact version and shown on the abstract page https://arxiv.org/abs/2608.20681v1. Full rights notice: licenses/arxiv-2608.20681-CC-BY-4.0.txt.\par\smallskip

\noindent\textit{Editorial scope.} Counted unit: the Lemma whose source label is \texttt{lem:sl2-nilpotent-separation} in the article (source0.tex lines 480--493, source label \texttt{lem:sl2-nilpotent-separation}), delivered together with the article's own complete original proof of it (source0.tex lines 495--571, one \texttt{proof} environment of 77 lines). No mathematics is written, repaired or re-derived: statement and proof are the article's own text line for line. Required context retained uncounted: none. Every definition, display and label of those lines is retained unchanged. Omitted from this package and not redistributed: 1--479, 494--494, 572--793. Nothing inside a retained excerpt is omitted or altered: every retained body is delivered exactly as the article's lines are written, so no omission receipt of any kind accompanies this package. Citations: the retained excerpts carry no citation command at all, so no bibliography entry of the article is transcribed and no reference is redistributed. Reference adaptation: none was needed and none was made; every reference inside the delivered unit resolves inside it, so no unresolved reference or citation is produced. Printed slips kept literal: no printed word, symbol, hypothesis or number of the counted statement or of its proof was altered. Layout and class: the article's own class is \texttt{amsart} with packages including fontenc, inputenc, lmodern, microtype, mathtools, amssymb, enumitem, xcolor, hyperref, cleveref; the wrapper is the lane's pinned \texttt{amsart} editorial wrapper at 10pt on A4, which restates the theorem-like environments the retained text needs and adds the article's own macro definitions that the retained text needs (\texttt{\textbackslash Ad}, \texttt{\textbackslash F}, \texttt{\textbackslash id}), with extra wrapper packages (bm, bbm, dsfont, xspace, cancel, mathtools). The delivered numbering is the unit's own. Assets: no retained span contains a figure environment, a graphics-inclusion command, a PDF-page inclusion, a table or a TikZ picture; no image file is copied into this package, the package declares no asset dependency and no image file is redistributed with it. Licence: version arXiv:2608.20681v1 of this article is distributed under the Creative Commons Attribution 4.0 International licence, as recorded in the arXiv licence metadata for this exact version and shown on the abstract page https://arxiv.org/abs/2608.20681v1. A full rights notice is delivered at licenses/arxiv-2608.20681-CC-BY-4.0.txt. Characters: every retained excerpt is pure ASCII, so the delivered engine, XeLaTeX with its default Latin Modern text font, reports no missing character.
\begin{lemma}\label{lem:sl2-nilpotent-separation}
Let $e,h,f$ be the standard generators of $\mathfrak{sl}_2$,
and $V$ a finite-dimensional $\mathfrak{sl}_2$-module.
If $S\in\operatorname{End}_{\mathfrak{sl}_2}(V)$ is semisimple,
then there exist semisimple elements $q_0,q_1\in\mathfrak{sl}_2$,
such that, for every operator $T\in\operatorname{End}_{\mathfrak{sl}_2}(V)$ with $[S,T]=0$,
\[
(S+e+T)\ker(S+q_0)\cap(S+e+T)\ker(S+q_1)\cap(S+e)\ker T= 0.
\]
In particular, when $T=0$, it follows that
\[
 (S+e)\ker(S+q_0)\cap(S+e)\ker(S+q_1)=0.
\]
\end{lemma}

\begin{proof}
Since \(S\) acts semisimply and commutes with \(\mathfrak{sl}_2\), put
\[
  V_\lambda=\ker(S-\lambda\id_V),
  \qquad
  M_{\lambda,m}=\operatorname{Hom}_{\mathfrak{sl}_2}\bigl(L(m),V_\lambda\bigr),
\]
where \(L(m)\) denotes the simple \(\mathfrak{sl}_2\)-module of highest
weight \(m\).  The canonical evaluation maps give the joint isotypic
decomposition
\[
  V\cong
  \bigoplus_{\lambda,m}M_{\lambda,m}\otimes L(m),
\]
where \(\mathfrak{sl}_2\) acts on the second factor and \(S\) acts by the scalar
\(\lambda\).

Choose \(t\in\F^\times\) such that
\[
  \lambda+tj\ne0
\]
whenever \(M_{\lambda,m}\ne0\) and \(j\ne0\) is an \(h\)-weight of
\(L(m)\).  If \(v_{k,0}\) is a nonzero zero-weight vector in \(L(2k)\),
then
\[
  \ker(S+th)=\bigoplus_{k\ge0} M_{0,2k}\otimes\F v_{k,0}.
\]
Set
\[
  q_0=th,\qquad
  q_1=t\Ad(\exp_{SL_2}(f))h= t(h+2f).
\]
Both \(q_0\) and \(q_1\) are semisimple, and
\[
  \ker(S+q_1)=\exp_{SL_2}(f)\ker(S+q_0).
\]

It follows that
\[
(S+e+T)\ker(S+q_0)\cap (S+e+T)\ker(S+q_1)\subseteq \bigoplus_{k\ge0} M_{0,2k}\otimes L(2k).
\]
For $k=0$, both $S$ and $e$ act as zero on
$M_{0,0}\otimes L(0)$, so the corresponding component of
$(S+e)\ker T$ is zero. If $y$ lies in the threefold intersection on
$M_{0,2k}\otimes L(2k)$, $k\ge 1$, then
\[
\begin{aligned}
 y&=Tm_0\otimes v_{k,0}+m_0\otimes ev_{k,0}\\
  &=Tm_1\otimes \exp_{SL_2}(f)v_{k,0}+m_1\otimes e\exp_{SL_2}(f)v_{k,0}, \quad m_0,m_1\in M_{0,2k},
\end{aligned}
\]
and
\[
 y\in (S+e)\ker T\subseteq\ker T.
\]
We have
\[
Ty= T^2m_0\otimes v_{k,0}+ Tm_0\otimes ev_{k,0}= 0.
\]
The linear independence of $v_{k,0}$ and $ev_{k,0}$ gives \(Tm_0=0\).
Similarly, \(Tm_1=0\). Therefore
\[
 y=m_0\otimes ev_{k,0}
  =m_1\otimes e\exp_{SL_2}(f)v_{k,0},
\]
which forces $y=0$ by the linear independence of $ev_{k,0}$ and $e\exp_{SL_2}(f)v_{k,0}$.
Thus the threefold intersection is zero on every joint summand on
which the first two spaces are supported, and it is therefore zero on $V$.

Finally, when $T=0$, one has $\ker T=V$, and the first two spaces
are contained in $(S+e)V$. Hence the threefold intersection reduces
to
\[
 (S+e)\ker(S+q_0)\cap(S+e)\ker(S+q_1),
\]
which proves the last assertion.
\end{proof}

\end{document}
